\documentclass[11pt,a4paper]{report}

% ----------------------------------------------------------------------
% packages
% ----------------------------------------------------------------------
\usepackage[T1]{fontenc}
\usepackage[utf8]{inputenc}
\usepackage{lmodern}
\usepackage[margin=2.6cm]{geometry}
\usepackage{microtype}
\usepackage{amsmath,amssymb,amsthm}
\usepackage{bm}
\usepackage{graphicx}
\usepackage{booktabs}
\usepackage{enumitem}
\usepackage{fancyvrb}
\usepackage{xcolor}
\usepackage{listings}
\usepackage[hidelinks]{hyperref}
\usepackage[nameinlink,capitalise]{cleveref}
\IfFileExists{generated/unicode.tex}{\input{generated/unicode.tex}}{}

% ----------------------------------------------------------------------
% notation
% ----------------------------------------------------------------------
\newcommand{\vect}[1]{\bm{#1}}
\newcommand{\X}{\vect{X}}
\newcommand{\Tr}{\operatorname{Tr}}
\newcommand{\ket}[1]{\lvert #1\rangle}
\newcommand{\bra}[1]{\langle #1\rvert}
\newcommand{\expval}[1]{\langle #1\rangle}
\newcommand{\comm}[2]{[#1,\,#2]}
\newcommand{\acomm}[2]{\{#1,\,#2\}}
\newcommand{\QL}{Q_L}
\newcommand{\QR}{Q_R}
\newcommand{\QV}{Q^V}
\newcommand{\QA}{Q^A}
\usepackage{url}
\DeclareUrlCommand\code{\urlstyle{tt}}   % code in text: breaks at . and _, no escaping needed

\theoremstyle{definition}
\newtheorem{definition}{Definition}[chapter]
\newtheorem{example}{Example}[chapter]
\theoremstyle{plain}
\newtheorem{theorem}{Theorem}[chapter]
\theoremstyle{remark}
\newtheorem*{remark}{Remark}

% ----------------------------------------------------------------------
% code listings
% ----------------------------------------------------------------------
\definecolor{codebg}{gray}{0.96}
\definecolor{codekw}{RGB}{0,82,147}
\definecolor{codecomment}{RGB}{110,110,110}
\definecolor{codestr}{RGB}{163,21,21}
\lstdefinestyle{python}{
  language=Python, basicstyle=\ttfamily\small, backgroundcolor=\color{codebg},
  keywordstyle=\color{codekw}\bfseries, commentstyle=\color{codecomment}\itshape,
  stringstyle=\color{codestr}, showstringspaces=false, columns=fullflexible,
  frame=none, xleftmargin=0.6em, framexleftmargin=0.6em, aboveskip=0.8em, belowskip=0.8em,
  breaklines=true, upquote=true}
\lstdefinestyle{shell}{basicstyle=\ttfamily\small, backgroundcolor=\color{codebg}, columns=fullflexible,
  xleftmargin=0.6em, framexleftmargin=0.6em, breaklines=true, upquote=true}
\lstnewenvironment{python}{\lstset{style=python}}{}
\lstnewenvironment{shell}{\lstset{style=shell}}{}

% ----------------------------------------------------------------------
% macros used by the generated API reference (tools/build_docs.py)
% ----------------------------------------------------------------------
\DefineVerbatimEnvironment{docverbatim}{Verbatim}{fontsize=\small,xleftmargin=1.2em,formatcom=\color{black!80}}
\newcommand{\docsection}[1]{\par\smallskip\noindent\textbf{#1}\par}
\newlist{docitems}{description}{1}
\setlist[docitems]{font=\normalfont,leftmargin=1.2em,labelindent=0pt,itemsep=1pt,style=nextline}
\newcommand{\docsig}[2]{\par\smallskip\noindent{\raggedright\ttfamily\small\textbf{#1}#2\par}\smallskip}
\newcommand{\docclass}[2]{\subsection*{\texttt{class #1}}\label{api:#2}%
  \addcontentsline{toc}{subsection}{\texttt{#1}}}
\newcommand{\docfunction}[2]{\subsection*{\texttt{#1}}\label{api:#2}}
\newcommand{\docmember}[2]{\par\medskip\noindent\textbf{\texttt{#1}}\enspace{\small\itshape #2}\par}
\newcommand{\docnote}[1]{\par\noindent{\small\itshape #1}\par}

% figures: generated per figure by tools/build_docs.py from figures/manifest.json
\newcommand{\missingfigure}[2]{\fbox{\parbox{0.8\linewidth}{\centering\small
  Figure \texttt{#1} is not available yet.\\ Create it with \texttt{#2}.}}}
\newcommand{\latsimfigure}[1]{\IfFileExists{generated/figures/#1.tex}{\input{generated/figures/#1.tex}}{%
  \begin{figure}[tbp]\centering\missingfigure{#1}{python tools/build\_docs.py}\caption{(missing)}\end{figure}}}

\setcounter{tocdepth}{1}
\numberwithin{equation}{chapter}

\title{\textbf{LatticeSim}\\[0.5em]
  \large Hamiltonian lattice simulations of bosonized (1+1)-dimensional gauge theories\\[0.3em]
  \normalsize From quantum mechanics to the search for bound states}
\author{LatticeSim documentation, release 0.1}
\date{\today}

\begin{document}
\maketitle
\tableofcontents

\chapter{Introduction}\label{chap:intro}

\section{What LatticeSim does}
LatticeSim simulates quantum field theories in one space and one time dimension on a spatial lattice, in the
Hamiltonian formulation. Its present scope is the class of theories that can be solved \emph{exactly} on the
lattice by linear algebra: \emph{massless fermions coupled to abelian gauge fields}, $U(1)^k$, with arbitrary
charges for the left- and right-moving fermions. After bosonization (\cref{chap:bosonization}) and elimination
of the gauge field through the Gauss law (\cref{chap:gauge}) such a theory becomes a \emph{quadratic}
Hamiltonian of coupled harmonic oscillators, and everything about it follows from one real symmetric matrix
(\cref{chap:gaussian}).

The goal of this release is a complete, tested path from the definition of a theory to the question
\begin{quote}
\emph{Which gauge-invariant particles does the theory contain, and which of them are bound states of the
elementary fermions?}
\end{quote}
The answer for the 3-4-5-0 chiral gauge theory, a standard test case, is worked out in \cref{chap:bound}:
the theory contains a massless gauge-invariant composite fermion $\psi_{3L}^2\psi_{4L}\psi_{5R}^{-2}$ and its
antiparticle, found both from the lattice Hamiltonian and from an independent analytic argument.

\section{How to read this document}
The prerequisites are quantum mechanics (harmonic oscillator, commutators, Heisenberg picture) and linear algebra.
Everything else is developed step by step:
\begin{description}[leftmargin=1.5em,style=nextline]
  \item[\Cref{chap:gaussian}] Quadratic Hamiltonians, normal modes, the Williamson normal form, Gaussian ground
    states and zero modes. This is the numerical core of the library.
  \item[\Cref{chap:lattice}] A free scalar field on a lattice: discretization, boundary conditions, the
    continuum limit, and the zero modes of a field on a ring.
  \item[\Cref{chap:bosonization}] Fermions in one dimension and their description by a compact boson; why the
    fermion numbers are windings and momenta of the boson.
  \item[\Cref{chap:gauge}] Gauge fields in one dimension: the Gauss law, the Schwinger model and its mass,
    correlation functions, screening and confinement of probe charges, the $\theta$ angle.
  \item[\Cref{chap:anomalies}] Gauge anomalies, the anomaly-cancellation condition, and which global symmetries
    survive the gauging.
  \item[\Cref{chap:bound}] Charge sectors, the bound-state search, and the 3-4-5-0 model as worked example,
    including a recipe for testing your own theory.
\end{description}
Each chapter states which part of the library implements it, and the figures are produced by the studies that
come with the library, so every number in this document can be reproduced.

\section{Quick start}
Installation (Python $\geq 3.10$) into a virtual environment, from the repository root:
\begin{shell}
python3 -m venv .venv
source .venv/bin/activate
pip install -e ".[dev]"
pytest                      # fast tests; "pytest -m slow" runs the expensive ones
\end{shell}
The bound-state search of \cref{chap:bound} in a few lines:
\begin{python}
from latsimpkg import ChiralTheory, find_bound_states

theory = ChiralTheory(QL=[3, 4], QR=[5, 0], g=0.3)   # charges of left and right movers
scan = find_bound_states(theory)

print(scan.agrees)                 # lattice and analytic prediction agree
for s in scan.bound_states:
    print(s.content, s.h, s.hb)    # psi_3L^-2 psi_4L^-1 psi_5R^2   0.5  0.0  (antiparticle)
                                   # psi_3L^2 psi_4L psi_5R^-2      0.5  0.0
\end{python}
Studies, figures and this document:
\begin{shell}
python -m studies               # run all studies, results -> results/
python -m studies.figures       # figures -> figures/ (only what is out of date)
python tools/build_docs.py      # this PDF -> docs/build/latticesim.pdf
\end{shell}

\section{Structure of the library}
The package \code{latsimpkg} is built in layers; every layer only uses the layers above it in this list.
\begin{center}
\small
\begin{tabular}{@{}lp{0.7\linewidth}@{}}
\toprule
module & content \\
\midrule
\code{conventions} & normalizations and sign conventions (\cref{app:conventions}) \\
\code{theory} & \code{ChiralTheory}: charges, anomaly check, pairing, meson masses \\
\code{lattice} & \code{Lattice}: number of sites, spacing, open chain or ring \\
\code{phasespace} & \code{PhaseSpace}: layout of the vector of coordinates and momenta \\
\code{quadratic} & \code{QuadraticHamiltonian}: normal modes, Williamson form, ground state, zero modes \\
\code{models} & \code{build_gauge_model}: theory + lattice $\to$ Hamiltonian \\
\code{analysis} & spectra, correlators, potentials, charge sectors, result types \\
\code{api} & \code{find_bound_states}, \code{spectrum}, \code{screening_potential} \\
\code{plotting} & figures of saved results \\
\bottomrule
\end{tabular}
\end{center}
The API reference (\cref{chap:api}) is generated from the docstrings of these modules.

\chapter{Quadratic Hamiltonians and Gaussian states}\label{chap:gaussian}

Every model in this library reduces to a system of coupled harmonic oscillators. This chapter develops the
complete theory of such systems: equations of motion, normal modes, the ground state, and the special role of
zero modes. It is implemented in \code{latsimpkg.quadratic.QuadraticHamiltonian}.

\section{Phase space and the symplectic matrix}
Consider $n$ degrees of freedom with coordinates $q_1,\dots,q_n$ and momenta $p_1,\dots,p_n$, collected in the
\emph{phase-space vector}
\begin{equation}
  \X = (q_1,\dots,q_n,\,p_1,\dots,p_n)^T .
\end{equation}
The canonical commutation relations $\comm{q_j}{p_k} = i\delta_{jk}$, $\comm{q_j}{q_k}=\comm{p_j}{p_k}=0$
($\hbar = 1$) take the compact form
\begin{equation}\label{eq:ccr}
  \comm{X_a}{X_b} = i J_{ab}, \qquad
  J = \begin{pmatrix} 0 & \mathbb{1}_n \\ -\mathbb{1}_n & 0 \end{pmatrix}.
\end{equation}
$J$ is the \emph{symplectic matrix}; it is real, antisymmetric, orthogonal, and $J^2 = -\mathbb{1}$.
In the library the layout of $\X$ is described by a \code{PhaseSpace} with named blocks (for a field: one block
of coordinates $\varphi_{i,j}$ per flavor $i$ and site $j$, and its conjugate momenta), and \code{PhaseSpace.J}
is the matrix \eqref{eq:ccr}.

\section{Quadratic Hamiltonians and equations of motion}
A quadratic Hamiltonian is
\begin{equation}\label{eq:quadratic-H}
  H = \tfrac12 \X^T M \X = \tfrac12 \sum_{a,b} M_{ab} X_a X_b ,
\end{equation}
with a real symmetric $2n\times 2n$ matrix $M$. Stability requires $M$ to be positive semidefinite,
$\X^T M \X \geq 0$, otherwise the energy is unbounded from below.

The Heisenberg equation $\dot X_a = i\comm{H}{X_a}$ follows from \eqref{eq:ccr} and
$\comm{X_bX_c}{X_a} = X_b\comm{X_c}{X_a} + \comm{X_b}{X_a}X_c$:
\begin{equation}
  \comm{H}{X_a} = \tfrac12 M_{bc}\bigl(i X_b J_{ca} + i J_{ba} X_c\bigr) = -i (JM\X)_a ,
\end{equation}
using $J^T = -J$ and the symmetry of $M$. Hence
\begin{equation}\label{eq:eom}
  \dot{\X} = J M \X , \qquad \X(t) = e^{JMt}\X(0).
\end{equation}
The operator equations are linear and identical to the classical Hamilton equations; everything about the
dynamics is encoded in the matrix $JM$.

\begin{example}[single oscillator]
For $H = \frac12(\omega^2 q^2 + p^2)$ we have $M = \operatorname{diag}(\omega^2, 1)$ and
$JM = \begin{pmatrix}0&1\\-\omega^2&0\end{pmatrix}$ with eigenvalues $\pm i\omega$.
\end{example}

\section{Normal modes and the Williamson normal form}
A linear change of variables $\X = S\vect{Y}$ preserves the commutation relations, $\comm{Y_a}{Y_b} = iJ_{ab}$,
if and only if
\begin{equation}
  S J S^T = J ,
\end{equation}
i.e.\ if $S$ is \emph{symplectic}. The quantum mechanics of quadratic Hamiltonians is therefore the linear
algebra of symplectic transformations. The central result is
\begin{theorem}[Williamson~\cite{Williamson1936}]\label{thm:williamson}
For every real symmetric positive definite $M$ there is a symplectic $S$ with
\begin{equation}
  S^T M S = \begin{pmatrix} \Omega & 0 \\ 0 & \Omega \end{pmatrix}, \qquad
  \Omega = \operatorname{diag}(\omega_1,\dots,\omega_n),\ \omega_k > 0 .
\end{equation}
\end{theorem}
In the new variables $\vect{Y} = (Q_1,\dots,Q_n,P_1,\dots,P_n)$ the Hamiltonian decouples into independent
oscillators,
\begin{equation}
  H = \sum_{k=1}^n \frac{\omega_k}{2}\bigl(Q_k^2 + P_k^2\bigr)
    = \sum_k \omega_k \bigl(a_k^\dagger a_k + \tfrac12\bigr),\qquad a_k = \tfrac{1}{\sqrt2}(Q_k + iP_k).
\end{equation}
The $\omega_k$ are the \emph{normal-mode frequencies}; the eigenvalues of $JM$ are $\pm i\omega_k$. The ground
state $\ket{0}$ is annihilated by all $a_k$ and has the \emph{zero-point energy}
\begin{equation}\label{eq:E0}
  E_0 = \frac12\sum_k \omega_k .
\end{equation}
Excited states are obtained by acting with $a_k^\dagger$; in field theory each $a_k^\dagger$ creates a particle
of energy $\omega_k$.

\paragraph{How the library computes it.} Diagonalizing the non-symmetric matrix $JM$ is numerically delicate.
Instead, $JM = M^{-1/2}\,(M^{1/2} J M^{1/2})\,M^{1/2}$ is similar to the antisymmetric matrix
$A = M^{1/2} J M^{1/2}$, so the frequencies are the positive eigenvalues of the \emph{Hermitian} matrix $iA$.
Its eigenvectors $v = (x + iy)/\sqrt2$ for eigenvalue $\omega$ satisfy $Ax = \omega y$, $Ay = -\omega x$. With the
orthogonal $O = (y_1,\dots,y_n,x_1,\dots,x_n)$ one has $O^T A O = \begin{pmatrix}0&\Omega\\-\Omega&0\end{pmatrix}$,
and
\begin{equation}
  S = M^{-1/2}\, O \begin{pmatrix}\Omega^{1/2}&0\\0&\Omega^{1/2}\end{pmatrix}
\end{equation}
is the Williamson transformation: $S^TMS = \operatorname{diag}(\Omega,\Omega)$ and $SJS^T = J$ (check it using
$M^{-1/2} J M^{-1/2} = -A^{-1}$). This is \code{QuadraticHamiltonian.frequencies} and \code{.williamson}.

\section{The ground state: Gaussian states and the covariance matrix}
The ground state of a quadratic Hamiltonian is a \emph{Gaussian state}: all its properties follow from the
first and second moments. By symmetry $\expval{\X} = 0$, and the second moments are collected in the
\emph{covariance matrix}
\begin{equation}
  \sigma_{ab} = \tfrac12 \expval{\acomm{X_a}{X_b}} .
\end{equation}
In normal-mode variables the oscillator ground state has $\expval{Q_k^2} = \expval{P_k^2} = \frac12$ and no
correlations, so $\sigma_Y = \frac12\mathbb{1}$, and transforming back,
\begin{equation}\label{eq:sigma}
  \sigma = \tfrac12 S S^T
  = \tfrac12 M^{-1/2}\,\bigl|M^{1/2} J M^{1/2}\bigr|\,M^{-1/2},
\end{equation}
where $|A| = \sqrt{A^TA}$. For a single oscillator this gives the familiar
$\expval{q^2} = 1/(2\omega)$ and $\expval{p^2} = \omega/2$. Any correlation function of linear operators
follows, e.g.\ $\expval{(\vect{u}^T\X)(\vect{v}^T\X)}_{\rm sym} = \vect{u}^T\sigma\vect{v}$; the electric-field
correlator of \cref{chap:gauge} is computed this way.

\paragraph{Uncertainty principle and purity.} For every quantum state
\begin{equation}\label{eq:uncertainty}
  \sigma + \tfrac{i}{2}J \;\geq\; 0
\end{equation}
(as a Hermitian matrix), the matrix form of the Heisenberg uncertainty relation~\cite{Weedbrook2012}. A Gaussian
state is \emph{pure} exactly if the smallest eigenvalue of $\sigma + \frac{i}{2}J$ is zero: all Williamson
eigenvalues of $\sigma$ equal $\frac12$. The tests check this for every ground state the library produces.

\section{Zero modes: free particles and conserved charges}\label{sec:zero-modes}
If $M$ has a kernel, \cref{thm:williamson} does not apply. The simplest example is a free particle,
$H = p^2/2m$: $M = \operatorname{diag}(0, 1/m)$, the coordinate $q$ is a \emph{flat direction}, and the momentum
$p$ is conserved. On a ring the uniform shift of a field is such a flat direction, and its conserved momentum
turns out to be a fermion number (\cref{chap:bosonization}); zero modes are therefore not a nuisance but the
place where the charges live.

\paragraph{Conserved linear charges.} A linear operator $\vect{w}^T\X$ is conserved if
$\frac{d}{dt}\vect{w}^T\X = \vect{w}^T J M \X = 0$ for all $\X$, i.e.
\begin{equation}
  M J \vect{w} = 0 \quad\Longleftrightarrow\quad \vect{w} \in J\ker M .
\end{equation}
With an orthonormal basis $K$ of $\ker M$ the conserved charges are the columns of $A = JK$.

\paragraph{Normal form with zero modes.} We require the kernel to be \emph{isotropic}, $K^TJK = 0$: no two flat
directions are conjugate to each other (a pair $q,p$ with $H$ independent of both would have no dynamics at all
and no ground state). Then every flat direction is a free particle, and the phase space splits into
\begin{itemize}
  \item the flat directions $K$ (positions of the free particles),
  \item their conjugates, measured by the conserved charges $\vect{c} = A^T\X$ (momenta of the free particles),
  \item the oscillator space $W = \{\X : K^T\X = 0,\ K^TJ\X = 0\}$.
\end{itemize}
On $W$ the matrix $M$ is positive definite and the Williamson construction applies with the symplectic form
$\Omega_W = B^TJB$ of a basis $B$ of $W$. Every $\X$ decomposes uniquely as
$\X = M^+ A G \vect{c} + B\vect{\xi} + K\vect{\kappa}$ with the pseudo-inverse $M^+$ and
$G = (A^T M^+ A)^{-1}$, and the energy separates exactly:
\begin{equation}\label{eq:free-normal-form}
  H = \tfrac12\, \vect{c}^T G\, \vect{c} \;+\; \tfrac12\, \vect{\xi}^T (B^TMB)\, \vect{\xi} .
\end{equation}
The first term is the kinetic energy of the free particles, the second a set of oscillators.

\paragraph{Lowest energy at fixed charges.} Because the two terms in \eqref{eq:free-normal-form} are independent,
the lowest energy in the sector with conserved charges $\vect{c}$ is $\frac12\vect{c}^TG\vect{c}$ above the
oscillator ground state. It equals the classical constrained minimum, which follows from a Lagrange multiplier
$\vect{\lambda}$:
\begin{equation}\label{eq:constrained-minimum}
  \min_{A^T\X = \vect{c}} \tfrac12 \X^T M \X : \quad M\X = A\vect{\lambda}
  \ \Rightarrow\ \X^* = M^+ A G \vect{c}, \quad E(\vect{c}) = \tfrac12\, \vect{c}^T G\, \vect{c} .
\end{equation}
This formula is the engine of the bound-state search (\cref{chap:bound}). In the library it is implemented by
the methods \code{free_energy} and \code{constrained_minimum} of \code{QuadraticHamiltonian}. The ground-state
covariance with zero
modes, \code{.covariance}, describes the oscillators at sharp conserved charges; the flat directions themselves
have infinite variance.

\section{In code}
\begin{python}
import numpy as np
from latsimpkg.quadratic import QuadraticHamiltonian

H = QuadraticHamiltonian(np.diag([1.7**2, 1.0]))   # one oscillator, omega = 1.7
H.frequencies             # [1.7]
H.covariance[0, 0]        # <q^2> = 1/(2 omega)
H.ground_state_energy     # 0.85

free = QuadraticHamiltonian(np.diag([0.0, 1.0]))   # free particle H = p^2/2
free.n_free, free.conserved_charges               # 1, A = J K: the charge c = -p
free.free_energy([2.0])                            # c = 2 (p = -2)  ->  E = 2
\end{python}
\code{QuadraticHamiltonian} checks its input: $M$ must be symmetric, positive semidefinite, finite, match the
phase space and have an isotropic kernel; otherwise it raises an error instead of returning wrong numbers. The
test suite verifies the construction on random symplectic matrices with known frequencies and free particles.

\chapter{A free scalar field on a lattice}\label{chap:lattice}

A quantum field in one spatial dimension is an infinite collection of coupled oscillators, one per point $x$.
Replacing the continuum by a lattice of $N$ sites turns it into exactly the finite system of
\cref{chap:gaussian}; this Hamiltonian lattice formulation goes back to Kogut and
Susskind~\cite{KogutSusskind1975}. This chapter introduces the discretization used throughout the library, the two boundary
conditions, and the continuum limit.

\section{The continuum theory}
A real scalar field $\varphi(x)$ with conjugate momentum $\Pi(x)$, $\comm{\varphi(x)}{\Pi(y)} = i\delta(x-y)$,
has the Hamiltonian
\begin{equation}\label{eq:H-continuum}
  H = \int dx\, \Bigl[\tfrac12 \Pi^2 + \tfrac12 (\partial_x\varphi)^2 + \tfrac12 \mu^2 \varphi^2\Bigr].
\end{equation}
Plane waves $e^{ikx}$ are normal modes with $\omega_k = \sqrt{\mu^2 + k^2}$: the field describes free particles of
mass $\mu$ (speed of light $c = 1$). For $\mu = 0$ the particles are massless and move at the speed of light.

\section{Discretization}
Put the field on sites $x_j = ja$, $j = 0,\dots,N-1$, with lattice spacing $a$, and define
\begin{equation}
  \varphi_j = \varphi(x_j), \qquad P_j = a\,\Pi(x_j) .
\end{equation}
The factor $a$ makes the lattice variables canonical: the continuum $\delta(x-y)$ becomes $\delta_{jk}/a$, so
$\comm{\varphi_j}{P_k} = i\delta_{jk}$. Replacing $\int dx \to a\sum_j$ and
$\partial_x\varphi \to (\varphi_{j+1}-\varphi_j)/a$ gives
\begin{equation}\label{eq:H-lattice}
  H = \sum_j \Bigl[\frac{P_j^2}{2a} + \frac{(\varphi_{j+1}-\varphi_j)^2}{2a}\Bigr]
     = \tfrac12 \X^T M \X, \qquad \X = (\varphi_0,\dots,\varphi_{N-1},P_0,\dots,P_{N-1}).
\end{equation}
The differences $\ell_j = \varphi_{j+1} - \varphi_j$ live on the \emph{links} between sites; the library
represents them by a matrix $B$ with $\ell = B\X$, so that the gradient energy is $\X^TB^TB\X/(2a)$.

\section{Open chain with Dirichlet boundary conditions}
On the open chain the field is fixed to zero just outside the lattice, $\varphi_{-1} = \varphi_N = 0$. There
are $N+1$ links and
\begin{equation}
  M = \frac1a\begin{pmatrix} K & 0 \\ 0 & \mathbb{1}\end{pmatrix},\qquad
  K = \begin{pmatrix} 2 & -1 & & \\ -1 & 2 & -1 & \\ & \ddots & \ddots & \ddots \\ & & -1 & 2\end{pmatrix}.
\end{equation}
The eigenvectors of $K$ are standing waves $\sin(k_n x_j)$ with
\begin{equation}\label{eq:dirichlet}
  k_n = \frac{\pi n}{L},\quad L = (N+1)a, \qquad
  \omega_n = \frac{2}{a}\sin\frac{k_n a}{2},\qquad n = 1,\dots,N .
\end{equation}
$L$ is the distance between the two fixed points. For $k_n a \ll 1$ the lattice dispersion reproduces the
continuum, $\omega_n = k_n\bigl(1 - (k_na)^2/24 + \dots\bigr)$: lattice artefacts are of order $(ka)^2$.
\Cref{fig:free_chain_spectrum} compares the numerical frequencies with \eqref{eq:dirichlet}.

\latsimfigure{free_chain_spectrum}

\paragraph{Ultraviolet behaviour.} The zero-point energy $E_0 = \frac12\sum_n\omega_n$ grows like $N/a$: it
is dominated by the shortest wavelengths and has no continuum limit. Physical statements always concern
\emph{differences} of energies (masses, excitation energies, potentials), which do have a limit.

\section{The ring and its zero modes}\label{sec:ring}
On a ring (periodic boundary conditions, $\varphi_N \equiv \varphi_0$, $L = Na$) the momenta are
$k_n = 2\pi n/L$, and the mode $n = 0$ is special: shifting the whole field, $\varphi_j \to \varphi_j + c$,
costs no energy. In the language of \cref{sec:zero-modes} the uniform shift is a flat direction, and the total
momentum
\begin{equation}
  P_{\rm tot} = \sum_j P_j
\end{equation}
is a conserved charge with kinetic energy $P_{\rm tot}^2/(2L)$ (each $P_j = P_{\rm tot}/N$ contributes
$P_j^2/2a$).

\paragraph{Compact fields and windings.} In the theories of this library the field is \emph{compact}: $\varphi$
and $\varphi + \sqrt\pi$ describe the same physical configuration (\cref{chap:bosonization} explains why). Then
the field may wind around the ring,
\begin{equation}
  \varphi(x + L) = \varphi(x) + \sqrt\pi\, n,\qquad n\in\mathbb{Z},
\end{equation}
and the winding number $n$ cannot change under continuous time evolution. The library writes such a field as a
periodic part plus a uniform slope, $\varphi_j = \tilde\varphi_j + W j/N$ with $W = \sqrt\pi n$ (in the
electric field the slope is centred, see \cref{chap:gauge}). Because the
periodic differences sum to zero, the gradient energy splits into $\sum_j(\tilde\varphi_{j+1}-\tilde\varphi_j)^2/2a
+ W^2/(2L)$. $W$ is treated as a canonical momentum, conjugate to a new coordinate $\bar\vartheta$ (the
\emph{dual zero mode}, \cref{chap:bosonization}): the phase space on the ring is
\begin{equation}
  \X = (\tilde\varphi_0,\dots,\tilde\varphi_{N-1},\ \bar\vartheta,\ P_0,\dots,P_{N-1},\ W),
\end{equation}
of dimension $2(N+1)$ per field. Both $P_{\rm tot}$ and $W$ are conserved for the free field, and the energy of a
configuration with given values is
\begin{equation}\label{eq:zero-mode-energy}
  E = \frac{P_{\rm tot}^2}{2L} + \frac{W^2}{2L} .
\end{equation}
With $P_{\rm tot} = \sqrt\pi m$ and $W = \sqrt\pi n$ this is $E = \pi(m^2 + n^2)/(2L)$: an energy proportional to
$1/L$, the hallmark of a massless particle in a box. \Cref{chap:bosonization} identifies $n$ and $m$ with fermion
numbers.

\section{In code}
\begin{python}
from latsimpkg import ChiralTheory, Lattice, build_gauge_model

free = ChiralTheory([0], [0], g=0.0)          # one neutral fermion = one free boson
open_chain = build_gauge_model(free, Lattice(50, a=0.1, boundary="open"))
open_chain.H.frequencies                        # (2/a) sin(k_n a/2)

ring = build_gauge_model(free, Lattice(20, boundary="ring"))
ring.H.n_free                                   # 2 conserved charges: P_tot and W
\end{python}

\chapter{Fermions in one dimension and bosonization}\label{chap:bosonization}

The theories of interest contain fermions, but the library simulates bosons. In one spatial dimension this is
not an approximation: a massless Dirac fermion \emph{is} a compact massless boson, a statement known as
bosonization~\cite{Coleman1975,Mandelstam1975,Giamarchi2003}. This chapter explains the part of the dictionary
the library uses and checks it on the spectrum of free fermions on a ring.

\section{Chiral fermions}
In one dimension a massless fermion field splits into a left-moving and a right-moving component,
$\psi = (\psi_L, \psi_R)$, with
\begin{equation}
  H = \int dx\,\Bigl[\psi_R^\dagger(-i\partial_x)\psi_R + \psi_L^\dagger(i\partial_x)\psi_L\Bigr].
\end{equation}
A right mover has energy $E = k$ for momentum $k$, a left mover $E = -k$; both move at the speed of light and
never scatter into each other without interactions. Each component has its own particle number,
$N_L = \int\psi_L^\dagger\psi_L$ and $N_R = \int\psi_R^\dagger\psi_R$ (measured relative to the filled Dirac sea).
A \emph{chiral} theory treats them differently, for example by giving them different gauge charges.

\paragraph{Spectrum on a ring.} On a ring of length $L$ (with antiperiodic boundary conditions for the fermions)
the single-particle momenta are $(r+\frac12)\,2\pi/L$. Adding $N_R$ right movers above the Fermi sea fills the
lowest levels and costs
\begin{equation}
  \frac{2\pi}{L}\sum_{r=0}^{N_R-1}\Bigl(r+\tfrac12\Bigr) = \frac{2\pi}{L}\,\frac{N_R^2}{2},
\end{equation}
the same for holes ($N_R<0$) and for left movers. The lowest energy at given fermion numbers is therefore
\begin{equation}\label{eq:fermion-sector}
  E(N_L, N_R) = \frac{2\pi}{L}\,(h + \bar h),\qquad h = \frac{N_L^2}{2},\quad \bar h = \frac{N_R^2}{2},
\end{equation}
and its momentum is $P = \frac{2\pi}{L}(\bar h - h)$ (left movers carry negative momentum). Above these states
there are particle-hole excitations with energies in multiples of $2\pi/L$. The numbers $h,\bar h$ are the
\emph{conformal weights} of the state; $\Delta = h + \bar h$ is its scaling dimension and $s = h - \bar h$ its
spin. A single fermion has $\Delta = |s| = \frac12$.

\section{The bosonization dictionary}
Introduce a scalar field $\varphi$ with momentum $\Pi$ such that the vector (total) and axial (left minus right)
fermion densities are
\begin{equation}\label{eq:dictionary}
  \rho_V = \psi_L^\dagger\psi_L + \psi_R^\dagger\psi_R = \frac{1}{\sqrt\pi}\,\partial_x\varphi,\qquad
  \rho_A = \psi_L^\dagger\psi_L - \psi_R^\dagger\psi_R = \frac{1}{\sqrt\pi}\,\Pi .
\end{equation}
The Hamiltonian of the free fermion becomes the free massless boson,
$H = \int\frac12[\Pi^2 + (\partial_x\varphi)^2]$, and the particle-hole excitations of the fermions are the
oscillator modes of the boson. It is convenient to define the \emph{dual field}
\begin{equation}
  \vartheta(x) = \int^x dy\,\Pi(y), \qquad \rho_A = \frac{1}{\sqrt\pi}\partial_x\vartheta ,
\end{equation}
so that $\varphi$ counts vector charge and $\vartheta$ axial charge in the same way.

\paragraph{Fermion numbers are zero modes.} Integrating \eqref{eq:dictionary} over the ring,
\begin{equation}
  n \equiv N_L + N_R = \frac{1}{\sqrt\pi}\bigl[\varphi(L) - \varphi(0)\bigr],\qquad
  m \equiv N_L - N_R = \frac{1}{\sqrt\pi}\int_0^L \Pi .
\end{equation}
The vector fermion number is the \emph{winding} of $\varphi$ and the axial number its \emph{total momentum}.
Since fermion numbers are integers, $\varphi$ must be compact with period $\sqrt\pi$ (\cref{sec:ring}), and $n$
and $m$ are integers with
\begin{equation}\label{eq:parity}
  n \equiv m \pmod 2,
\end{equation}
because $N_L = (n+m)/2$ and $N_R = (n-m)/2$ are integers. This parity rule is what makes the compact boson
fermionic: $(n,m) = (1,0)$ would be half a left and half a right fermion.

\paragraph{Check.} The zero-mode energy of the boson \eqref{eq:zero-mode-energy} with $W = \sqrt\pi n$,
$P_{\rm tot} = \sqrt\pi m$ is
\begin{equation}
  E = \frac{\pi(n^2+m^2)}{2L} = \frac{\pi}{L}\bigl(N_L^2 + N_R^2\bigr) = \frac{2\pi}{L}\Bigl(\frac{N_L^2}{2}
  + \frac{N_R^2}{2}\Bigr),
\end{equation}
exactly the fermionic result \eqref{eq:fermion-sector}. The lattice momentum $P = -\frac1a\sum_j P_j\ell_j$
gives
\begin{equation}
  s = -\frac{PL}{2\pi} = \frac{nm}{2} = \frac{N_L^2 - N_R^2}{2} = h - \bar h .
\end{equation}
The library reproduces both exactly on any ring
(\code{analysis.sectors.sector_energy}; tested for several $(N_L,N_R)$). Note the sign convention: left movers have
negative momentum and positive spin, $s = h - \bar h$.

\section{Many fermions and their charges}
A theory contains left movers $L_1,\dots,L_{n_L}$ and right movers $R_1,\dots,R_{n_R}$, with charges
$Q_{L,i}$, $Q_{R,j}$ under a $U(1)$ gauge field (for $U(1)^k$: one row of charges per gauge field). To bosonize,
the library pairs $L_i$ with $R_i$ into a Dirac fermion and describes each pair by one compact boson
$\varphi_i$. The gauge current $\sum_i (Q_{L,i}\rho_{L,i} + Q_{R,i}\rho_{R,i})$ becomes
\begin{equation}\label{eq:va-charges}
  \sum_i \bigl(Q^V_i \rho_{V,i} + Q^A_i \rho_{A,i}\bigr),\qquad
  Q^V = \frac{Q_L + Q_R}{2},\quad Q^A = \frac{Q_L - Q_R}{2}.
\end{equation}
The pairing is bookkeeping: physical results must not depend on it, and the tests check this explicitly for the
two possible pairings of the 3-4-5-0 model. If $n_L \neq n_R$, the missing partners are added as neutral fermions
(charge $0$), which are free and decoupled (\code{ChiralTheory.pairing}).

\begin{python}
from latsimpkg import ChiralTheory

t = ChiralTheory(QL=[3, 4], QR=[5, 0], g=0.3)
p = t.pairing()
p.QV, p.QA          # [[4, 2]], [[-1, 2]]
t.pairing(order=[1, 0]).QV   # the other pairing: [[1.5, 4.5]]
\end{python}

\chapter{Gauge fields in one dimension}\label{chap:gauge}

In one spatial dimension the electromagnetic field has no transverse polarizations: there are no photons. The
electric field is completely fixed by the charges through the Gauss law, and eliminating it leaves an
interaction between the charges. For bosonized fermions this interaction is quadratic, so the gauge theory stays
a system of coupled oscillators. This chapter derives the lattice Hamiltonian used by
\code{models.build_gauge_model} and its main physical consequences.

\section{Electric field and Gauss law}
In the temporal gauge $A_0 = 0$ the only field is $A_1(x)$ with conjugate momentum $-E(x)$, and the Hamiltonian
of the gauge field is $\frac12\int E^2$ (there is no magnetic field in one dimension). Physical states satisfy
the Gauss law
\begin{equation}\label{eq:gauss}
  \partial_x E = g\,\rho ,
\end{equation}
where $\rho$ is the charge density of the matter and $g$ the coupling. With the bosonized densities
\eqref{eq:dictionary} and charges \eqref{eq:va-charges},
\begin{equation}
  g\rho = \frac{g}{\sqrt\pi}\sum_i\bigl(Q^V_i\,\partial_x\varphi_i + Q^A_i\,\partial_x\vartheta_i\bigr)
  \quad\Rightarrow\quad
  E(x) = \frac{g}{\sqrt\pi}\sum_i\bigl(Q^V_i\varphi_i(x) + Q^A_i\vartheta_i(x)\bigr) + e(x) .
\end{equation}
The integration constant $e(x)$ is a c-number background: a constant $e = g\theta/2\pi$ (the $\theta$ angle) or
the field of static external charges (\cref{sec:probes}). The total Hamiltonian is
\begin{equation}\label{eq:H-gauge}
  H = \int dx\,\Bigl[\tfrac12\sum_i\bigl(\Pi_i^2 + (\partial_x\varphi_i)^2\bigr) + \tfrac12 E^2\Bigr].
\end{equation}
With several gauge fields $U(1)^k$ there is one $E_a$ with coupling $g_a$ and charges $Q^V_a, Q^A_a$ per gauge
field, and the electric energy is $\frac12\sum_a\int E_a^2$.

\paragraph{On the lattice.} With $\vartheta_{i,j} = \sum_{k\leq j}P_{i,k}$ (the discrete version of
$\int^x\Pi$) the electric field on site $j$ is a linear function of $\X$,
\begin{equation}
  E_j = (C\X)_j + e_j,\qquad
  E_j = \frac{g}{\sqrt\pi}\sum_i\Bigl(Q^V_i\varphi_{i,j} + Q^A_i\sum_{k\leq j}P_{i,k}\Bigr) + e_j ,
\end{equation}
and the electric energy $\frac a2\sum_jE_j^2$ adds $aC^TC$ to the matrix $M$ (plus a linear term and a constant
if $e\neq0$). On the ring the zero modes enter as well:
$\varphi_{i,j} = \tilde\varphi_{i,j} + W_ij/N$ and $\vartheta_{i,j} = \bar\vartheta_i + \sum_{k\leq j}P_{i,k}$.

\section{The Schwinger model}
The Schwinger model~\cite{Schwinger1962,LowensteinSwieca1971} is a single Dirac fermion with charge $1$
($Q_L = Q_R = 1$, so $Q^V = 1$, $Q^A = 0$). Then $E = (g/\sqrt\pi)\varphi$ and \eqref{eq:H-gauge} becomes
\begin{equation}
  H = \int dx\,\Bigl[\tfrac12\Pi^2 + \tfrac12(\partial_x\varphi)^2 + \frac{g^2}{2\pi}\varphi^2\Bigr]:
\end{equation}
a free boson of mass
\begin{equation}\label{eq:schwinger-mass}
  \mu = \frac{g}{\sqrt\pi} .
\end{equation}
The massless fermions have disappeared from the spectrum; the only particle is a massive neutral boson, the
``meson''. Its mass is generated non-perturbatively by the gauge field.

\paragraph{General charges.} For one gauge field with charges $Q_L,Q_R$ (anomaly free, \cref{chap:anomalies})
exactly one combination of the bosons becomes massive, with
\begin{equation}\label{eq:meson-mass}
  \mu^2 = \frac{g^2}{\pi}\sum_i\bigl[(Q^V_i)^2 + (Q^A_i)^2\bigr] = \frac{g^2}{2\pi}\sum\bigl(Q_L^2 + Q_R^2\bigr),
\end{equation}
and the remaining $F-1$ combinations stay massless. For $U(1)^k$ the masses are the square roots of the
eigenvalues of $\mu^2_{ab} = (g_ag_b/2\pi)(Q_{L,a}\cdot Q_{L,b} + Q_{R,a}\cdot Q_{R,b})$
(\code{ChiralTheory.meson_masses}).

\section{Spectra on the open chain}
On the open chain the predicted spectrum is one branch $\sqrt{\mu^2 + k_n^2}$ and $F-1$ massless branches $k_n$,
with the Dirichlet momenta \eqref{eq:dirichlet}. For vector charges the electric field acts locally on $\varphi$
and the prediction is exact on the lattice. Axial charges act on $\vartheta$, a sum over all sites to the left,
whose boundary conditions differ from Dirichlet: the momenta become $k_n = (n+\delta)\pi/L$ with a phase shift
$\delta$, and the prediction holds only in the bulk. \Cref{fig:spectra_open_chain,fig:finite_size_scaling} show
both effects; the deviation for axial charges falls like $1/N$.

\latsimfigure{spectra_open_chain}
\latsimfigure{finite_size_scaling}

\paragraph{Continuum limit.} At fixed physical length $L$ the lowest frequency of the Schwinger model approaches
$\sqrt{\mu^2 + (\pi/L)^2}$ as $a\to0$, with the $(ka)^2$ lattice correction of \cref{chap:lattice}
(\cref{fig:continuum_limit}).

\latsimfigure{continuum_limit}

\section{Measuring the mass: the electric-field correlator}
The electric field is proportional to the massive boson, so its correlations measure $\mu$. For a free boson of
mass $\mu$ the equal-time correlator in the ground state is
\begin{equation}
  \expval{\Phi(x)\Phi(0)} = \int\frac{dk}{2\pi}\,\frac{e^{ikx}}{2\sqrt{k^2+\mu^2}} = \frac{1}{2\pi}K_0(\mu|x|),
\end{equation}
with the modified Bessel function $K_0$, which decays like $e^{-\mu|x|}$. The library computes
$\expval{E_jE_k} = (C\sigma C^T)_{jk}$ from the covariance matrix \eqref{eq:sigma} and fits $A K_0(\mu r) + B$
in a window away from the boundaries (\code{analysis.correlators.meson_mass_from_E} with an explicit
\code{FitWindow}). The offset $B$ is a uniform fluctuation of the background field caused by the open boundary,
non-zero for axial charges. The fitted masses agree with \eqref{eq:meson-mass} to about $10^{-3}$
(\cref{fig:correlator_fits}).

\latsimfigure{correlator_fits}

\section{Probe charges: screening and confinement}\label{sec:probes}
Place static external charges $+q$ at $x=0$ and $-q$ at $x = d$. By the Gauss law they create the background
$e(x) = gq$ for $0<x<d$. The ground-state energy shift is the potential $V(d)$ between them. For the Schwinger
model the Hamiltonian is $\int[\frac12\Pi^2 + \frac12(\partial\varphi)^2 + \frac12(\mu\varphi + e)^2]$, and
minimizing over $\varphi$ in Fourier space gives
\begin{equation}
  V(d) = \frac12\int\frac{dk}{2\pi}\,|e_k|^2\frac{k^2}{k^2+\mu^2}
       = g^2q^2\int\frac{dk}{2\pi}\frac{1 - \cos kd}{k^2+\mu^2}
       = \frac{g^2q^2}{2\mu}\bigl(1 - e^{-\mu d}\bigr).
\end{equation}
The potential saturates: the massless fermions \emph{screen} any charge, even a fractional one~\cite{CJS1975}.
With $\mu = g/\sqrt\pi$ the prefactor is $g^2q^2/2\mu = \frac{\sqrt\pi}{2}gq^2$, which is exactly the flat-space
result quoted in~\cite{BharadwajIsenKang2026}, eq.~(3.88); the lattice reproduces it to $5\times10^{-4}$.
Without charged matter ($\mu = 0$) the same formula gives $V(d) = g^2q^2d/2$: a constant electric field between
the charges, linear \emph{confinement} with string tension $g^2q^2/2$. On the lattice the potential is the
background energy $\min_{\X}[\frac12\X^TM\X + \vect{b}^T\X] + \text{const} = \text{const} - \frac12\vect{b}^TM^+\vect{b}$
(\code{api.screening_potential}); \cref{fig:probe_potentials} shows both regimes.

\latsimfigure{probe_potentials}

\section{The theta angle}
A constant background field $e = g\theta/2\pi$ is the $\theta$ angle of the theory. With massless fermions it is
unphysical: in the Schwinger model it is removed by the shift $\varphi\to\varphi - \sqrt\pi\,\theta/2\pi$, which
corresponds to a chiral rotation of the fermions. On the ring this shift is a zero mode and costs nothing, so the
background energy vanishes (numerically $10^{-15}$). On the open chain the Dirichlet boundary fixes $\varphi$ at the
ends and a small boundary energy remains. A physical $\theta$ dependence requires massive fermions, which are
beyond the Gaussian methods of this release.

\section{In code}
\begin{python}
from latsimpkg import schwinger, ChiralTheory, spectrum, screening_potential

sp = spectrum(schwinger(g=0.04), N=300)          # open chain
sp.max_deviation                                 # ~1e-14: exact for vector charges
pot = screening_potential(schwinger(g=0.1), N=400, distances=range(2, 120, 6), q=0.5)
pot.meta["regime"], pot.max_rel_deviation        # 'screening', ~5e-4
\end{python}

\chapter{Anomalies and global symmetries}\label{chap:anomalies}

Not every assignment of charges defines a consistent gauge theory. A chiral theory can suffer from a
\emph{gauge anomaly}: a symmetry of the classical theory that is necessarily broken by quantization. For a gauge
symmetry this is fatal. On the lattice the anomaly shows up in a simple and concrete way: the electric field
does not commute with itself.

\section{The anomaly on the lattice}
The Gauss law says that the gauge transformations are generated by $G(x) = \partial_xE - g\rho$. For a consistent
theory these generators must commute; equivalently, the electric field determined by the Gauss law must satisfy
$\comm{E(x)}{E(y)} = 0$. Using $\comm{\varphi_j}{\vartheta_k} = i\,\mathbb{1}(j\leq k)$, which follows from
$\vartheta_k = \sum_{l\leq k}P_l$, one finds for one gauge field
\begin{equation}\label{eq:EE}
  \comm{E_j}{E_k} = i\,\frac{g^2}{\pi}\,(Q^V\!\cdot Q^A)\,\operatorname{sign}(k-j),\qquad
  Q^V\!\cdot Q^A = \frac14\Bigl(\sum Q_L^2 - \sum Q_R^2\Bigr).
\end{equation}
The electric field commutes with itself if and only if
\begin{equation}\label{eq:anomaly-free}
  \sum_i Q_{L,i}^2 = \sum_j Q_{R,j}^2 ,
\end{equation}
the anomaly-cancellation condition of two-dimensional chiral gauge theories~\cite{Bertlmann1996}. The coefficient
$\sum Q_L^2 - \sum Q_R^2$ is the anomaly: it counts the mismatch of the charged left- and right-movers.
\Cref{fig:anomalous_commutator} shows \eqref{eq:EE} for several anomalous charge assignments.

\latsimfigure{anomalous_commutator}

\begin{example}
The Schwinger model ($1^2 = 1^2$) and its axial version are anomaly free. The 3-4-5-0 model has
$3^2 + 4^2 = 25 = 5^2 + 0^2$: a genuinely chiral but consistent theory. Charges $Q_L = (2,1)$, $Q_R = (0,1)$
give $5 - 1 = 4 \neq 0$ and are rejected by \code{ChiralTheory} unless \code{allow_anomalous=True}.
\end{example}

\paragraph{Several gauge fields.} For $U(1)^k$ every pair of gauge fields must be anomaly free,
\begin{equation}
  \mathcal{A}_{ab} = Q_{L,a}\cdot Q_{L,b} - Q_{R,a}\cdot Q_{R,b} = 0,
\end{equation}
including the \emph{mixed} anomalies $a\neq b$ (\code{ChiralTheory.anomaly_matrix}). In vector-axial form
$\mathcal{A}_{ab} = 2(Q^V_a\cdot Q^A_b + Q^A_a\cdot Q^V_b)$.

\paragraph{A lattice subtlety for $k>1$.} On the lattice
$\comm{E_{a,j}}{E_{b,k}} = i\frac{g_ag_b}{\pi}\bigl(x_{ab}\mathbb{1}(j\leq k) - x_{ba}\mathbb{1}(k\leq j)\bigr)$
with $x_{ab} = Q^V_a\cdot Q^A_b$. Anomaly freedom only gives $x_{ab} + x_{ba} = 0$, so a constant commutator
$i x_{ab}$ between \emph{different} gauge fields can remain. Since
$x_{ab} = (Q_{R,a}\cdot Q_{L,b} - Q_{L,a}\cdot Q_{R,b})/4$ depends on how left and right movers are paired, it is
an artefact of the regularization, absent for pairings with $Q_LQ_R^T$ symmetric.
\code{build_gauge_model} searches for such a pairing automatically.

\section{Which global symmetries survive}
Each fermion number is a global $U(1)$ symmetry of the free theory: $n_L + n_R$ of them, with charges
$t = (t_L, t_R)$. After gauging, a global symmetry survives only if it has no \emph{mixed} anomaly with the
gauge field,
\begin{equation}\label{eq:mixed}
  t_L\cdot Q_{L,a} = t_R\cdot Q_{R,a}\qquad\text{for all } a .
\end{equation}
The others are broken by the anomaly, just as the axial charge of the Schwinger model is not conserved. For an
anomaly-free theory the gauge charges themselves satisfy \eqref{eq:mixed}. The number of surviving symmetries is
$n_L + n_R - \operatorname{rank}[Q_L, -Q_R]$ (\code{ChiralTheory.global_symmetries}): one for the Schwinger model
(the gauge charge itself), three for the 3-4-5-0 model.

\paragraph{The lattice finds them by itself.} On the ring the conserved charges are the kernel directions of
$MJ$ (\cref{sec:zero-modes}). In terms of the fermion numbers, a conserved combination $t_n\cdot n + t_m\cdot m$
corresponds to $t_L\cdot N_L + t_R\cdot N_R$ with $t_n = (t_L+t_R)/2$, $t_m = (t_L-t_R)/2$, and \eqref{eq:mixed}
becomes $Q^A_a\cdot t_n + Q^V_a\cdot t_m = 0$. The tests verify that the conserved charges found numerically from
the lattice Hamiltonian span exactly the space of \eqref{eq:mixed}: the anomaly structure is not put in by hand, it
emerges from $M$.

\chapter{Charge sectors and the search for bound states}\label{chap:bound}

We now have all ingredients to answer the question of the introduction: which gauge-invariant particles does a
theory contain? After the gauge field has given a mass to some combinations of the bosons, the remaining
massless bosons are compact, and their charge sectors are the low-energy particles. Some of them are the original
fermions, others are \emph{composites}: gauge-invariant bound states of several fermions. The search is
implemented in \code{find_bound_states}.

\section{The idea}
A state with fermion numbers $N_L = (N_{L,1},\dots)$ and $N_R = (N_{R,1},\dots)$ carries gauge charge
\begin{equation}
  q_a = Q_{L,a}\cdot N_L + Q_{R,a}\cdot N_R = Q^V_a\cdot n + Q^A_a\cdot m ,
\end{equation}
with $n = N_L + N_R$ and $m = N_L - N_R$. On a ring the Gauss law $\oint\partial_xE = gq = 0$ allows only
\emph{neutral} sectors, $q_a = 0$ for all $a$. Neutral combinations of charged fermions are candidates for bound
states. Whether such a sector really contains a light particle is decided by its energy: on a ring of length $L$
a massless particle has energy $E\propto1/L$, so
\begin{equation}
  \Delta = \frac{EL}{2\pi} = h + \bar h
\end{equation}
stays finite as $L\to\infty$ for a massless sector, while a sector that only contains massive excitations does not
produce a $1/L$ level at all.

\section{The lattice computation}
On the ring every flavor has the zero modes $W_i = \sqrt\pi n_i$ and $P_{{\rm tot},i} = \sqrt\pi m_i$
(\cref{sec:ring}). The gauge field couples some combinations of them to the massive meson, so they are no longer
conserved; the conserved charges are the kernel directions of $MJ$, i.e.\ the anomaly-free global symmetries of
\cref{chap:anomalies}. The algorithm is:
\begin{enumerate}
  \item Build the ring Hamiltonian $M$ (\code{build_gauge_model} with \code{boundary="ring"}) and its conserved
    charges $A$ (\cref{sec:zero-modes}).
  \item For every neutral $(N_L,N_R)$ with $|N|\leq n_{\max}$: the configuration $\X_v$ with
    $W_i = \sqrt\pi n_i$, $P_{i,j} = \sqrt\pi m_i/N$ has conserved charges $\vect{c} = A^T\X_v$.
  \item Vectors with the same $\vect{c}$ are the same sector (they differ by excitations of the massive meson); keep
    the representative with the smallest $|N_L|^2 + |N_R|^2$.
  \item The lowest energy in the sector is $E = \frac12\vect{c}^TG\vect{c}$ \eqref{eq:constrained-minimum}: only
    the conserved charges are fixed, all other directions relax into the massive mode. The spin follows from the
    lattice momentum of the minimizing configuration, $s = -PL/2\pi$.
  \item Classify: $\Delta = |s| = \frac12$ is a single massless fermion; it is \emph{elementary} if the
    representative is one of the original fermions and a \emph{bound state} otherwise. All other sectors are
    multi-particle states.
\end{enumerate}
Because the zero-mode energies are exact on the lattice (\cref{chap:bosonization}), $\Delta$ is independent of $N$
and of $g>0$: the result is not an extrapolation.

\section{The analytic cross-check: the charge lattice}
The same answer follows from a short argument. The gauge fields give a mass to the combinations
$Q_{L,a}\cdot\varphi_L$ and $Q_{R,a}\cdot\varphi_R$ of the left- and right-moving bosons. What remains massless is
the part of $(N_L, N_R)$ orthogonal to the row spaces of $Q_L$ and $Q_R$:
\begin{equation}\label{eq:charge-lattice}
  h = \tfrac12\bigl|N_L^\perp\bigr|^2,\qquad \bar h = \tfrac12\bigl|N_R^\perp\bigr|^2,\qquad
  s = h - \bar h = \tfrac12\bigl(|N_L|^2 - |N_R|^2\bigr),
\end{equation}
for neutral $(N_L,N_R)$; the last equality uses neutrality. These vectors form the \emph{charge lattice} of the
low-energy theory (\code{analysis.charge_lattice}). \code{find_bound_states} computes both and reports whether
they agree (\code{SectorScan.agrees}).

\section{Worked example: the 3-4-5-0 model}
The 3-4-5-0 model~\cite{WangWen2019,Tong2022} has two left movers with charges $3,4$ and two right movers with
charges $5,0$ under one $U(1)$.
\begin{description}[leftmargin=1.5em,style=nextline]
  \item[Anomaly] $3^2 + 4^2 = 25 = 5^2 + 0^2$: consistent, although no fermion is vector-like.
  \item[Pairing] $(3,5)$ and $(4,0)$: $Q^V = (4,2)$, $Q^A = (-1,2)$; indeed $Q^V\cdot Q^A = 0$.
  \item[Meson] $\mu^2 = \frac{g^2}{2\pi}(9 + 16 + 25 + 0) = \frac{25g^2}{\pi}$, one massive and one massless boson.
  \item[Global symmetries] four fermion numbers, one condition \eqref{eq:mixed}: three conserved charges, one of
    them the gauge charge.
\end{description}
The search over $|N|\leq3$, $\Delta\leq2$ finds twelve sectors, in exact agreement with the charge lattice
(\cref{fig:sectors_3450}):
\begin{center}
\begin{tabular}{@{}lllll@{}}
\toprule
$N_L$ & $N_R$ & $(h,\bar h)$ & kind & content \\
\midrule
$(2,1)$ & $(-2,0)$ & $(\frac12,0)$ & bound state & $\psi_{3L}^2\psi_{4L}\psi_{5R}^{-2}$ \\
$(-2,-1)$ & $(2,0)$ & $(\frac12,0)$ & bound state & its antiparticle \\
$(0,0)$ & $(0,\pm1)$ & $(0,\frac12)$ & elementary & $\psi_{0R}^{\pm1}$ \\
four sectors & & $(\frac12,\frac12)$ & multi-particle & bound state $\times\ \psi_{0R}^{\pm1}$ \\
four sectors & & $(2,0)$, $(0,2)$ & multi-particle & two identical fermions \\
\bottomrule
\end{tabular}
\end{center}
The bound state is neutral, $2\cdot3 + 1\cdot4 - 2\cdot5 = 0$, and left-moving with $h = \frac12$: from
\eqref{eq:charge-lattice}, $N_L^\perp$ is the component of $(2,1)$ orthogonal to $(3,4)$, of length
$|4\cdot2 - 3\cdot1|/5 = 1$, and $N_R^\perp = 0$. It is a massless gauge-invariant composite fermion. Together with
the neutral right mover $\psi_{0R}$ it forms the massless low-energy content of the theory: one left- and one
right-moving fermion, central charges $(c_L, c_R) = (1,1)$ instead of $(2,2)$ before gauging. Counting particle and
antiparticle, the theory has two bound states.

\latsimfigure{sectors_3450}

\paragraph{Comparison with the literature.} The same four fermions carry a second anomaly-free $U(1)$,
$H = (5,0\,|\,3,-4)$, without mixed anomaly with $G = (3,4\,|\,5,0)$~\cite{Tong2022}; the library finds $H$ among the
surviving global symmetries of the gauged 3-4-5-0 model. Tong gauges $H$ and finds at low energies a single massless
Dirac fermion, the elementary left mover $\Psi = \psi_2$ and the gauge-invariant composite right mover
$\tilde\Psi = \psi_1^2(\tilde\psi_1^\dagger)^2\tilde\psi_2$, with $G$-charges $+4$ and $-4$~\cite{Tong2022}.
\code{find_bound_states(ChiralTheory([5, 0], [3, -4]))} reproduces exactly this: $\psi_2$ is elementary with
$(h,\bar h) = (\frac12,0)$, and the only bound states are $\psi_1^{\pm2}\tilde\psi_1^{\mp2}\tilde\psi_2^{\pm1}$ with
$(h,\bar h) = (0,\frac12)$ and $G$-charge $\mp4$. Gauging $G$ instead, as in this chapter, is the same theory up to
parity ($L\leftrightarrow R$) and charge conjugation of one fermion, and the sectors coincide with $h$ and $\bar h$
exchanged (\code{tests/test_literature.py}). Wang and Wen~\cite{WangWen2019} show that the 3-4-5-0 theory can be
realized as the low-energy limit of a local lattice model with an on-site $U(1)$ symmetry.

\paragraph{Volume dependence.} \Cref{fig:sector_volume_dependence} shows $EL/2\pi$ against the ring size. The
neutral massless sectors stay exactly at $h+\bar h$. Charged sectors are forbidden by the Gauss law on the ring; if
one computes them anyway, the seam of the ring acts as an external counter-charge that the massless matter screens
(\cref{sec:probes}), so $E$ stays finite and $EL$ grows: no massless particle.

\latsimfigure{sector_volume_dependence}

\section{Testing your own theory}
\begin{python}
from latsimpkg import ChiralTheory, find_bound_states
from latsimpkg.plotting import plot_sector_diagram, save_figure

theory = ChiralTheory(QL=[3, 4], QR=[5], g=0.3)   # rejects anomalous charges
print(theory.anomaly_matrix, theory.meson_masses, theory.global_symmetries().shape[1])

scan = find_bound_states(theory, N=40, nmax=3, Delta_max=2.0)
assert scan.agrees                                 # lattice = charge lattice
for s in scan.bound_states:
    print(s.content, s.h, s.hb)
save_figure(plot_sector_diagram(scan), "my_theory")
scan.save("my_theory.npz")                         # reload with SectorScan.load
\end{python}
In this example the right mover with charge $0$ is missing; the pairing adds it as a neutral spectator
(\code{spectator_R1}), and the bound states are those of the 3-4-5-0 model.

\paragraph{Limitations.}
\begin{itemize}
  \item The coupling must resolve the meson in double precision, $\mu L\gtrsim10^{-2}$ and $\mu a\lesssim10$
    (\cref{sec:numerics}); outside this window \code{find_bound_states} raises an error instead of returning sectors.
    The result itself does not depend on $g$: for 3-4-5-0 it is identical from $g = 10^{-4}$ to $10^2$.
  \item Only massless fermions with abelian gauge fields: the Hamiltonian must be quadratic. Massive fermions or
    four-fermion interactions add cosine terms to the bosonized Hamiltonian and require non-Gaussian methods.
  \item The search finds massless sectors. Massive bound states would be additional poles in correlation functions;
    in the Gaussian theories of this release the only massive particles are the free mesons.
  \item \code{nmax} limits which composites can be found: a bound state needs a representative with
    $|N_{L,i}|,|N_{R,i}|\leq n_{\max}$.
  \item The classification assumes a fermionic massless sector ($h\in k^2/2$), as for 3-4-5-0; for other
    theories inspect $(h,\bar h)$ directly.
\end{itemize}


\appendix
\chapter{Conventions and checks}\label{app:conventions}

\section{Conventions}
All conventions are defined once in \code{latsimpkg.conventions}.
\begin{center}
\begin{tabular}{@{}ll@{}}
\toprule
quantity & convention \\
\midrule
units & $\hbar = c = 1$ \\
phase space & $\X = (q_1..q_n, p_1..p_n)$, $\comm{X_a}{X_b} = iJ_{ab}$, $J = \begin{pmatrix}0&\mathbb1\\-\mathbb1&0\end{pmatrix}$ \\
Hamiltonian & $H = \frac12\X^TM\X\ (+\vect{b}^T\X + \text{const with background})$ \\
lattice variables & $\varphi_j = \varphi(ja)$, $P_j = a\Pi(ja)$, $\comm{\varphi_j}{P_k} = i\delta_{jk}$ \\
length & open chain $L = (N+1)a$, ring $L = Na$ \\
dual field & $\vartheta_j = \sum_{k\leq j}P_k$ (+ $\bar\vartheta$ on the ring) \\
bosonization & $\rho_V = \partial_x\varphi/\sqrt\pi$, $\rho_A = \rho_L - \rho_R = \Pi/\sqrt\pi$, $\varphi\sim\varphi+\sqrt\pi$ \\
fermion numbers & $n = N_L + N_R$ (winding), $m = N_L - N_R$ (momentum), $n\equiv m \bmod 2$ \\
charges & $Q^V = (Q_L+Q_R)/2$, $Q^A = (Q_L-Q_R)/2$ \\
Gauss law & $E_j - E_{j-1} = g\,\rho_j$ (+ probe charges), background $e = g\theta/2\pi$ \\
spin & $s = -PL/2\pi = h - \bar h$, $P = -\frac1a\sum_jP_j\ell_j$; left movers have $s>0$ \\
\bottomrule
\end{tabular}
\end{center}

\section{What the test suite checks}
The library is tested against analytic results (\code{pytest}; \code{pytest -m slow} for the expensive tests):
\begin{itemize}
  \item linear algebra: oscillator moments, Williamson form on random symplectic matrices with known frequencies and
    free particles, symplecticity, purity of ground states, constrained minima, input validation;
  \item free field: Dirichlet spectrum to $10^{-12}$, zero-point energy, ring sectors
    $\Delta = (N_L^2+N_R^2)/2$ and $s = (N_L^2-N_R^2)/2$;
  \item gauge theories: Schwinger spectrum exact, meson masses, $a^2$ continuum limit, anomalous commutator
    \eqref{eq:EE}, correlator fit, screening and confinement potentials, $\theta$ on the ring;
  \item anomalies and symmetries: anomaly matrices, mixed anomalies, number of conserved charges, numerical
    conserved charges $=$ analytic global symmetries;
  \item bound states: 3-4-5-0 sectors and bound states, lattice $=$ charge lattice, independence of $g$ and $N$,
    $U(1)^2$ examples;
  \item invariances: pairing, flavor permutation, charge conjugation;
  \item infrastructure: saving and loading of all results, plotting without global state, incremental figures.
\end{itemize}

\chapter{API reference}\label{chap:api}
This chapter is generated from the docstrings of the library (\texttt{tools/build\_docs.py}). The docstrings keep
their formulas in plain ASCII so that they stay readable in an editor; the typeset mathematics behind them is in
the preceding chapters.
\IfFileExists{generated/api.tex}{\input{generated/api.tex}}{\par\emph{Run \texttt{python tools/build\_docs.py} to generate it.}}

\begin{thebibliography}{99}
\addcontentsline{toc}{chapter}{References}

\bibitem{Williamson1936}
J.~Williamson, \emph{On the algebraic problem concerning the normal forms of linear dynamical systems},
Am.\ J.\ Math.\ \textbf{58}, 141--163 (1936).

\bibitem{Weedbrook2012}
C.~Weedbrook, S.~Pirandola, R.~Garc\'ia-Patr\'on, N.~J.~Cerf, T.~C.~Ralph, J.~H.~Shapiro, S.~Lloyd,
\emph{Gaussian quantum information}, Rev.\ Mod.\ Phys.\ \textbf{84}, 621--669 (2012), arXiv:1110.3234.

\bibitem{Schwinger1962}
J.~Schwinger, \emph{Gauge invariance and mass.~II}, Phys.\ Rev.\ \textbf{128}, 2425--2429 (1962).

\bibitem{LowensteinSwieca1971}
J.~H.~Lowenstein and J.~A.~Swieca, \emph{Quantum electrodynamics in two dimensions},
Ann.\ Phys.\ \textbf{68}, 172--195 (1971).

\bibitem{CJS1975}
S.~Coleman, R.~Jackiw, L.~Susskind, \emph{Charge shielding and quark confinement in the massive Schwinger model},
Ann.\ Phys.\ \textbf{93}, 267 (1975).

\bibitem{Coleman1975}
S.~Coleman, \emph{Quantum sine-Gordon equation as the massive Thirring model}, Phys.\ Rev.\ D \textbf{11}, 2088--2097 (1975).

\bibitem{Mandelstam1975}
S.~Mandelstam, \emph{Soliton operators for the quantized sine-Gordon equation}, Phys.\ Rev.\ D \textbf{11}, 3026 (1975).

\bibitem{KogutSusskind1975}
J.~Kogut and L.~Susskind, \emph{Hamiltonian formulation of Wilson's lattice gauge theories},
Phys.\ Rev.\ D \textbf{11}, 395 (1975).

\bibitem{Giamarchi2003}
T.~Giamarchi, \emph{Quantum Physics in One Dimension}, Oxford University Press (2003).

\bibitem{Bertlmann1996}
R.~A.~Bertlmann, \emph{Anomalies in Quantum Field Theory}, Oxford University Press (1996).

\bibitem{WangWen2019}
J.~Wang and X.-G.~Wen, \emph{Solution to the 1+1 dimensional gauged chiral fermion problem},
Phys.\ Rev.\ D \textbf{99}, 111501(R) (2019), arXiv:1807.05998.

\bibitem{Tong2022}
D.~Tong, \emph{Comments on symmetric mass generation in 2d and 4d}, JHEP \textbf{07}, 001 (2022), arXiv:2104.03997.

\bibitem{BharadwajIsenKang2026}
S.~Bharadwaj, J.~Isen, Z.-B.~Kang, \emph{Confinement versus screening in the Schwinger model on AdS$_2$ from
bosonization and tensor networks}, arXiv:2607.19468 (2026).

\end{thebibliography}

\end{document}
